\documentclass[11pt, oneside]{article} 
\usepackage{geometry}
\geometry{letterpaper} 
\usepackage{graphicx}
	
\usepackage{amssymb}
\usepackage{amsmath}
\usepackage{parskip}
\usepackage{color}
\usepackage{hyperref}

\graphicspath{{/Users/telliott/Dropbox/Github-math/figures/}}
% \begin{center} \includegraphics [scale=0.4] {gauss3.png} \end{center}

\title{Computations}
\date{}

\begin{document}
\maketitle
\Large

%[my-super-duper-separator]

In this chapter we will continue with the theme of the \hyperref[sec:law_of_cosines]{\textbf{Law of cosines}}.  There is a general theorem and two specific cases.  I think it's perhaps easier if we start with the general case.

\subsection*{Stewart's theorem}

\label{sec:Stewarts_theorem}

\begin{center} \includegraphics [scale=0.5] {bisector3b.png} \end{center}

We have a triangle with sides $a$, $b$ and $c$ and a line segment connecting one vertex to the opposite side, dividing it into segments $x$ and $y$.  In the general case these two segments are not equal ($x \ne y$), and the vertex angle is not bisected ($\angle s \ne \angle t$).  The two angles at the base are $\theta$ and $\theta'$.  We use the law of cosines to write:

\emph{Proof}.

\[ a^2 = e^2 + x^2 - 2ex \cos \theta \]
\[ b^2 = e^2 + y^2 - 2ey \cos \theta' \]

Since $\theta$ and $\theta'$ are supplementary angles, $\cos \theta = - \cos \theta'$ so
\[ b^2 = e^2 + y^2 + 2ey \cos \theta \]

We divide the first equation by $x$ and the second by $y$ and add:
\[ \frac{a^2}{x} + \frac{b^2}{y} = \frac{e^2 + x^2}{x} + \frac{e^2 + y^2}{y} \]

Get rid of the fractions:
\[ a^2y + b^2x = e^2 (x+y) + yx^2 + xy^2 \]
\[ = (e^2 + xy) (x + y) \]

$\square$

This is called Stewart's theorem.

\subsection*{Apollonius' theorem}

\label{sec:Apollonius_theorem}

The first special case is where the line segment divides the base in half, so that $x = y$.  This is called Apollonius' theorem.  

\emph{Proof}.

We had

\[ a^2y + b^2x = (e^2 + xy) (x + y) \]

Substituting $m = x = y$
\[ a^2m + b^2m = (e^2 + mn)(2m) \]
\[ a^2 + b^2 = 2(e^2 + m^2) \]
\[ \frac{a^2 + b^2}{2} = e^2 + m^2 \]

$\square$

The average of the squares of the sides is equal to the median squared plus the half-base squared.

\subsection*{review of angle bisector theorem}

Let's go back to the generalized case of the angle bisector theorem, which we looked at  (\hyperref[sec:generalized_angle_bisector_theorem]{\textbf{here}}).

\begin{center} \includegraphics [scale=0.5] {bisector1.png} \end{center}

\[ \frac{x}{a} = \frac{y}{b} \]

As we turn the base $c$, but keep the angle bisection relationship, the ratio stays the same.

\emph{Proof}.

For a simple proof, consider the areas of the two smaller triangles. Twice the area of the top one is $a \cos s \cdot e$ (where $e$ is the length of the bisector).  Twice the area of the bottom one is $b$ times the same factor.  So the ratio of areas is $a/b$.

But if we take side $c = x + y$ as the base, then twice the area of the two triangles is the same altitude times either $x$ or $y$, so the ratio of the two areas is also $x/y$.

\[ \frac{x}{y} = \frac{a}{b} \]

$\square$

A somewhat different proof relies on the law of sines.  We have that 
\[ \frac{x}{\sin s} = \frac{a}{\sin \theta} \]
\[ \frac{y}{\sin t} = \frac{b}{\sin \theta'} \] 

\begin{center} \includegraphics [scale=0.5] {bisector3b.png} \end{center}

But $s = t$ so $\sin s = \sin t$ and also $\theta$ and $\theta'$ are supplementary so $\sin \theta = \sin \theta'$.  We obtain
\[ \frac{x}{a} = \frac{\sin s}{\sin \theta} = \frac{\sin t}{\sin \theta'} = \frac{y}{b} \]

$\square$

\subsection*{bisected angle and Stewart's theorem}

We use the result for the bisected angle to simplify Stewart's theorem for that case.

\begin{center} \includegraphics [scale=0.5] {bisector1.png} \end{center}

\emph{Proof}.

For a bisected angle, our basic relationship is 
\[ \frac{x}{a} = \frac{y}{b} \]

which can be rewritten in two ways as
\[ x = \frac{ay}{b} \]
\[ y = \frac{bx}{a} \]

and we can use this to simplify Stewart's theorem for the special case. 
\[ a^2y + b^2x = (e^2 + xy) \cdot (x + y) \]
\[ = (e^2 + xy) \cdot (\frac{bx}{a} + \frac{ay}{b}) \]
\[ = (e^2 + xy) \cdot (\frac{bx}{a} + \frac{ay}{b}) \cdot ab \cdot \frac{1}{ab} \]

The middle terms on the right-hand side cancel the left-hand side
\[ a^2y + b^2x = (\frac{bx}{a} + \frac{ay}{b}) \cdot ab \]

which leaves
\[ ab = e^2 + xy \]
\[ e^2 = ab - xy \]

$\square$

which was certainly not obvious to me.

\subsection*{median in terms of side lengths}

This problem is Hopkins 966.

\begin{center} \includegraphics [scale=0.4] {Hopkins_966.png} \end{center}

We just use Apollonius' theorem.

\[ a^2 + b^2 = 2(e^2 + m^2) \]

We have to change symbols.  We used $e$ for the line segment dropping from the vertex, but he has $m$.  If $c$ is the length of the base, then

\[ a^2 + b^2 = 2(m^2 + (c/2)^2) \]
\[ 2(a^2 + b^2) = 4m^2 + c^2) \]

\[ 4m^2 = 2(a^2 + b^2) - c^2 \]
\[ m = \frac{1}{2} \ \sqrt{2(a^2 + b^2) - c^2} \]

\subsection*{Hopkins 967}

\begin{center} \includegraphics [scale=0.5] {Hopkins_967.png} \end{center}

I thought, how hard can this be?  If you know Stewart's theorem, it's not hard.

First of all, we notice that $\triangle ACD \sim \triangle BDE$, because $\angle A = \angle E$ (equal arcs), and there is a shared vertical angle.  It turns out that doesn't help.

Their solution starts with the statement.

\[ k^2 = ab - AD \cdot DB \]

(switching from $e$ in our notation to $k$).

By this time we recognize this as the simplified form of Stewart's theorem for the special case of the bisected angle.  Leave that for the moment.  

Recall the result from the angle bisector theorem:

\[ \frac{AD}{b} = \frac{BD}{a} \]
\[ \frac{BD}{AD} = \frac{a}{b} \]
\[ \frac{BD + AD}{AD} = \frac{a + b}{b} \]
\[ \frac{c}{a + b} = \frac{AD}{b} = \frac{BD}{a} \]

So then
\[ AD = \frac{bc}{a + b} \]
\[ BD = \frac{ac}{a + b} \]

Finally, we can use the original expression to write:
\[ k^2 = ab -  \frac{bc}{a + b} \cdot \frac{ac}{a + b} \]

which can be manipulated a bit, without making it any simpler.  
\[ k^2 = ab -  \frac{bc}{a + b} \cdot \frac{ac}{a + b} \]
\[ = ab \ [1 - \frac{c^2}{(a + b)^2}] \]

We have an expression for the magnitude of the angle bisector in terms of the sides of the triangle.  It is symmetric in $a$ and $b$, which we should expect.


\end{document}